\section{齐性空间与截面表示}
\label{sec:ghs-homogeneous-space}

前面的 Lie 群表示由矩阵、根和权来描述。若表示空间由几何上的函数组成，表示的重数还受稳定子约束：同一个群可以作用在不同的空间上，而这些空间未必包含相同的不可约表示。球面说明这一点，旋转不仅作用在坐标上，也通过坐标变换作用在波函数上。

在球面上，场的角模态常由 $SO(d)$ 的不可约表示组织。
一般流形可能没有非平凡的等距群：自旋描述局部正交标架下的张量表示，
角动量描述全局等距群对态空间的作用；只有足够对称的几何才直接联系两者。

\begin{definition}[等距群的函数表示与齐性空间]
设 $(\Sigma,h)$ 是无边界 Riemann 流形，$\dd\mu_h$ 是其体积测度，$\Delta_h=\operatorname{div}_h\operatorname{grad}_h$ 是非正号的标量 Laplace 算子。
等距群 $G=\operatorname{Isom}(\Sigma,h)$ 由保持度规的微分同胚组成，函数空间采用
\[
 \langle f_1,f_2\rangle=\int_\Sigma\overline{f_1(x)}f_2(x)\,\dd\mu_h(x).
\]
若 $G$ 在 $\Sigma$ 上传递作用，选定 $o\in\Sigma$ 后其闭稳定子为
$H=\{g\in G:g\cdot o=o\}$，轨道映射给出 $\Sigma\simeq G/H$。
这里 $G/H$ 是齐性空间；$H$ 不必正规，因而一般没有商群乘法。
\end{definition}

令 $G=\operatorname{Isom}(\Sigma,h)$，其元素满足 $g^*h=h$。在 $L^2(\Sigma)$ 上定义
\begin{equation}\label{eq:ghs-isometry-unitary}
 (U_gf)(x)=f(g^{-1}x).
\end{equation}
因为群元 $g$ 保持体积测度，对任意函数 $f_1,f_2$ 作变量替换便有
$\langle U_gf_1,U_gf_2\rangle=\langle f_1,f_2\rangle$；故 $U$ 是幺正表示。与此同时 $\Delta_hU_g=U_g\Delta_h$，于是拉普拉斯算子的谱子空间在 $G$ 下不变。

\begin{proposition}[几何对称性导致谱的表示分解]
\label{prop:ghs-laplacian-isometry}
若 $\Sigma$ 为紧致流形，则每个 $-\Delta_h$ 的有限维特征空间是 $G$ 的有限维幺正表示；它可以分解为 $G$ 的不可约表示直和。若 $G$ 平凡，仍有拉普拉斯谱，却没有由 $G$ 赋予的角动量标签。
\end{proposition}
\begin{proof}
对任何 $f$ 有 $\Delta_h(f\circ g^{-1})=(\Delta_hf)\circ g^{-1}$，这是 Levi--Civita 联络的等距自然性。若 $-\Delta_hf=\lambda f$，立即有 $-\Delta_hU_gf=\lambda U_gf$。紧致性使每个谱特征空间有限维；其上的幺正表示由正交补不变性而完全可约。
\end{proof}


在传递作用的情形，场的角依赖可由 $G/H$ 上的函数表示理解。轨道--稳定子定理描述点的轨道，函数空间上的表示则把同一个构造用于全部角模式。

紧致群的 Haar 测度以下归一化为 $\int_G\dd g=1$。
正交性是前文 Schur 正交关系的 Haar 积分形式；矩阵系数的完备性由 Peter--Weyl 定理给出。

\begin{theorem}[紧致群正则表示的 Peter--Weyl 完备性]
不可约幺正表示的矩阵系数
$\sqrt{\dim V_\lambda}\,D^{(\lambda)}_{ij}(g)$ 构成 $L^2(G)$ 的正交完备基。
在左、右平移作用下，
\[
 L^2(G)\simeq\widehat{\bigoplus}_\lambda V_\lambda\otimes V_\lambda^*.
\]
\end{theorem}

这个紧致群完备性定理是以下证明使用的前提。
几何空间不是整个 $G$ 时，还需取稳定子所允许的矩阵系数；这一步可以直接计算。

\begin{theorem}[齐性空间的 Peter--Weyl 分解]
\label{thm:ghs-peter-weyl-homogeneous}
设 $G$ 为紧致群，$H\subset G$ 为闭子群，$V_\lambda$ 遍历 $G$ 的有限维不可约复表示。则
\begin{equation}\label{eq:ghs-homogeneous-peter-weyl}
 L^2(G/H)\simeq\widehat{\bigoplus}_{\lambda}\,V_\lambda\otimes(V_\lambda^*)^H,
\end{equation}
每个 $V_\lambda$ 出现的重数为 $\dim V_\lambda^H$。这里右上角 $H$ 表示不变子空间，帽号表示 Hilbert 空间的正交完备化。
\end{theorem}
\begin{proof}
首先把 $L^2(G/H)$ 识别为 $L^2(G)$ 中满足 $f(gh)=f(g)$ 的右 $H$ 不变函数。Peter--Weyl 定理给出
$L^2(G)\simeq\widehat\bigoplus_\lambda V_\lambda\otimes V_\lambda^*$，左 $G$ 只作用于第一因子，而右 $H$ 只作用于第二因子。逐个表示取右 $H$ 不变部分，得到 \cref{eq:ghs-homogeneous-peter-weyl}。因为 $V_\lambda$ 幺正，$\dim(V_\lambda^*)^H=\dim V_\lambda^H$，故重数公式成立。
\end{proof}

这个重数公式比直接记忆“球谐有多少个分量”更一般。若 $\Sigma=G/H$，只要知道紧致群 $G$ 的表示和稳定子 $H$ 的不变向量，就能判定哪些角模态可能出现。并且不需要 $\Sigma$ 是标准球面。

\begin{figure}[htbp]
\centering
% Local TikZ geometry; styles inherited exclusively from the shared preamble.
\begin{tikzfig}[scale=.94]
\draw[thick] (0,0) circle (1.65);
\draw[gray!70,dashed] (-1.65,0) arc[start angle=180,end angle=360,x radius=1.65,y radius=.44];
\draw[thick] (-1.65,0) arc[start angle=180,end angle=0,x radius=1.65,y radius=.44];
\draw[gray] (0,-2)--(0,2);
\fill (0,1.65) circle (1.7pt) node[above left] {$p$};
\draw[lecturearrow] (.24,1.93) arc[start angle=150,end angle=30,x radius=.46,y radius=.2];
\node[lecturelabel] at (1.33,2.15) {$H=SO(2)$};
\fill (1.65,0) circle (1.7pt);
\draw[lecturearrow] (1.35,.27) arc[start angle=28,end angle=133,x radius=1.4,y radius=.37];
\node[lecturelabel] at (-2.7,-.78) {$G\cdot q$};
\node[lecturelabel] at (0,-2.25) {$S^2\simeq SO(3)/SO(2)$};
\end{tikzfig}

\caption{球面的群论构造：$SO(3)$ 移动一点形成轨道 $S^2$，保持北极 $p$ 不变的子群是 $SO(2)$；因此 $S^2\simeq SO(3)/SO(2)$。}\label{fig:ghs-group-orbit-stabilizer}
\end{figure}

\cref{fig:ghs-group-orbit-stabilizer} 中轨道对应点的变化，稳定子对应对同一点不产生位移的旋转。先区分两者，才有办法判断 $L^2(G/H)$ 中哪些表示真正出现。


\begin{example}[球谐的表示重数]
把球面写成 $SO(3)/SO(2)$，由稳定子权零条件求 $L^2(S^2)$ 中各转动群表示的重数。

\textbf{解。} 令 $G=SO(3)$、$H=SO(2)$，其中 $H$ 是绕第三轴的转动。
自旋 $\ell$ 的表示在这个子群下分解为权 $m=-\ell,\ldots,\ell$，
权零向量恰有一个。因此 $L^2(S^2)$ 中每个整数 $\ell$ 表示出现一次，
其维数是 $2\ell+1$。半整数自旋不是 $SO(3)$ 的单值表示，不能作为标量球谐出现。
同一个 $\ell$ 内的 $m$ 简并由旋转群决定；不同 $\ell$ 的特征值则需指定具体算子。
\end{example}

函数的球谐可由齐性空间 $G/H$ 的标量函数分解得到；高自旋场却取值于表示纤维。这时不能再只保留 $H$ 的不变向量，而必须问：哪些 $G$ 不可约表示在限制到 $H$ 后包含给定的纤维表示？这个问题把群论中的表示限制与张量谐波统一起来。


\begin{definition}[关联丛与等变截面]
设 $G$ 紧致、$H\subset G$ 闭、$V_\rho$ 是 $H$ 的有限维幺正表示。齐性空间上的关联丛为
\[
 E_\rho=G\times_HV_\rho\longrightarrow G/H.
\]
它的平方可积截面等价于满足 $F(gh)=\rho(h^{-1})F(g)$ 的可测向量值函数 $F:G\to V_\rho$。群 $G$ 通过左平移作用在这些截面上。
\end{definition}

\begin{theorem}[张量截面的表示重数]
\label{thm:ghs-induced-bundle-decomp}
按 $G$ 不可约表示 $V_\lambda$ 分解，
\begin{equation}\label{eq:ghs-vectorbundle-peter-weyl}
 L^2(G\times_HV_\rho)\simeq
 \widehat{\bigoplus}_{\lambda}\;
 V_\lambda\otimes\operatorname{Hom}_H(V_\lambda,V_\rho).
\end{equation}
因此 $V_\lambda$ 的重数是 $\dim\operatorname{Hom}_H(V_\lambda,V_\rho)$。若 $V_\rho=\CC$ 为平凡表示，退化为 \cref{eq:ghs-homogeneous-peter-weyl}。
\end{theorem}
\begin{proof}
将向量值平方可积函数空间写成 $L^2(G)\otimes V_\rho$，其中 $H$ 在第一个因子的右正则表示与第二因子的纤维表示上作对角作用。利用
$L^2(G)\simeq\widehat\bigoplus_\lambda V_\lambda\otimes V_\lambda^*$，取满足右 $H$ 等变条件的子空间。
对有限维表示而言，$(V_\lambda^*\otimes V_\rho)^H$ 与 $\operatorname{Hom}_H(V_\lambda,V_\rho)$ 自然同构：一个张量 $\sum_i\ell_i\otimes w_i$ 对应线性映射 $v\mapsto\sum_i\ell_i(v)w_i$，右 $H$ 不变性等价于该映射与 $H$ 作用交换。于是得到分解与重数公式。
\end{proof}

\begin{definition}[球面一形式、Hodge 星与正 Laplace 算子]
在有定向的单位球面 $S^2$ 上，$\dd$ 是外微分，$*$ 把一形式转成相对于度规与定向旋转九十度的一形式。
其内积是 $\langle\alpha,\beta\rangle=\int_{S^2}\overline\alpha\wedge*\beta$。
令 $\delta=\dd^\dagger$，正定 Hodge Laplace 算子为
$\Delta_{\rm H}=\dd\delta+\delta\dd$。标量球谐满足
$\Delta_{\rm H}Y_{\ell m}=\ell(\ell+1)Y_{\ell m}$ 且范数为一。
\end{definition}

取 $\Sigma=S^2=SO(3)/SO(2)$。标量球谐 $Y_{\ell m}$ 的 $\ell=0,1,2,\dots$ 各出现一次。要在切丛上构造矢量模态，最直接的做法是用几何算子作用于标量球谐。对 $\ell\geq1$，定义
\begin{equation}\label{eq:ghs-vector-harmonics-sphere}
 E_{\ell m}=\frac{\dd Y_{\ell m}}{\sqrt{\ell(\ell+1)}},\qquad
 B_{\ell m}=\frac{*\dd Y_{\ell m}}{\sqrt{\ell(\ell+1)}}.
\end{equation}
这里 $E/B$ 表示两类正交的一形式谐波，不表示电磁场强度。

\begin{proposition}[球面矢量谐波的两类旋转多重态]

在单位球面 $S^2$ 上，$E_{\ell m}$ 与 $B_{\ell m}$ 都是 Hodge Laplacian 的本征一形式，本征值为 $\ell(\ell+1)$；对每个 $\ell\geq1$，共有 $2(2\ell+1)$ 个模式。
\end{proposition}
\begin{proof}
由 $\Delta_{\rm H}\dd=\dd\Delta_{\rm H}$ 和 $\Delta_{\rm H}*=*\Delta_{\rm H}$，直接从
$\Delta_{\rm H}Y_{\ell m}=\ell(\ell+1)Y_{\ell m}$ 得到相同的一形式本征值。积分分部给出
\[
 \norm{\dd Y_{\ell m}}^2
 =\langle Y_{\ell m},\delta\dd Y_{\ell m}\rangle
 =\ell(\ell+1),
\]
故归一化系数正确。$E$ 与 $B$ 的内积为
$\int_{S^2}\dd\overline Y\wedge\dd Y'$ 的常数倍，它是全微分，闭流形上积分为零。$H^1(S^2)=0$，没有额外调和一形式；Hodge 分解说明全部非零一形式都由恰当形式和余恰当形式生成，故模式数及完备性成立。
\end{proof}

\begin{example}[从稳定子权计算两类矢量模态]
将复化球面切纤维分解为稳定子的权 $+1,-1$，求各角动量下矢量模态的重数，并判断 $\ell=0$ 的情形。

\textbf{解。} 复化切纤维在 $SO(2)$ 下分解成权 $+1$ 与 $-1$ 两个一维表示。
$SO(3)$ 的自旋 $\ell$ 表示限制到 $SO(2)$ 后，每个 $m=-\ell,\ldots,\ell$ 恰出现一次。
因此对于 $\ell\geq1$，$\operatorname{Hom}_{SO(2)}(V_\ell,V_\rho)$ 的维数为二；
$\ell=0$ 不含 $\pm1$，故重数为零。总模式数是 $2\dim V_\ell=2(2\ell+1)$，
与 $E_{\ell m},B_{\ell m}$ 的几何构造相符。
\end{example}

这说明矢量球谐的两重类型来自微分形式的 Hodge 分解，标记 $\ell,m$ 则来自 $SO(3)$ 群表示。若将球面变成一般二维紧致流形，Hodge 分解仍成立，但 $SO(3)$ 的 $\ell,m$ 标签可能完全消失；若流形亏格为 $g$，还出现 $2g$ 个调和一形式零模。


\begin{proposition}[重数空间中的动力学混合]
设 $A$ 为有界算子，或为其谱投影均与 $G$ 对易的自伴算子。其在 \cref{eq:ghs-vectorbundle-peter-weyl} 的每个不可约分量上采取 $\id_{V_\lambda}\otimes A_\lambda$ 的形式。这里 $A_\lambda$ 作用在重数空间 $\operatorname{Hom}_H(V_\lambda,V_\rho)$ 上。重数为一时由 Schur 引理，$A$ 在该块上是标量；重数大于一时，它仍可混合多个等价 $G$ 表示，必须进一步对角化重数矩阵。
\end{proposition}
\begin{proof}
在两个不同不可约表示之间，算子矩阵元是零交织算符；同一个不可约表示内，Schur 引理使表示因子上的交织算符成为恒等算符的倍数。剩下的全部矩阵元素因而只作用在重数空间。对自伴情形先把此结论用于有界谱投影，再由谱演算恢复算子及其定义域。
\end{proof}

因此“对称性使算子对角化”并不总是“一块只有一个本征值”：真正由群决定的是不可约块，具体本征值与块内重数混合仍由动力学决定。这个细节在自旋更高、纤维表示更复杂时尤为重要。


在齐性空间上，群表示既描述函数空间，也描述张量丛的截面。将球谐的构造重新写成表示分解，就能同时确定模态的自旋、简并度与可允许的角向微分。
